%%%%%
%
%% Dans ce fichier, les �ventuels accents sont cod�s
%% en ISO-8859-15 (aka latin9)
%% La ligne suivante donne trois exemples 
%% e accent aigu : �, e accent grave : �, e-dans-l'o : �
\section{Position du soleil (depuis un emplacement mobile)}

\itemstyle{Op\'erateur~:} \operfmt{pos\_sol\_mobile\_loa}
\index{pos\_sol\_mobile\_loa}

%%%%%OBJET%%%%
\itemstyle{Objet~:}
Calcul des angles donnant la direction du soleil depuis
un endroit sur terre.

%%%%%ARGUMENTS%%%%
\bigskip
%% Dans ce fichier, les �ventuels accents sont cod�s
%% en ISO-8859-15 (aka latin9)
%% La ligne suivante donne trois exemples 
%% e accent aigu : �, e accent grave : �, e-dans-l'o : �
%%%%%ENTREES%%%%
\begin{tabarg}
\multicolumn{3}{l}{\titretab {Entr\'ees}}\\
\hline
$JDA$ & Coeff., \texttt{int} & Num�ro du jour dans l'ann�e (de 1 � 366)\\
$Time$ & Vecteur, \texttt{double} & Heure UTC (sec. depuis 0h00 du jour pass� en premier argument)\\
$Lon$ & Vecteur, \texttt{float} & Longitude (degree)\\
$Lat$ & Vecteur, \texttt{float} & Latitude (degree)\\
\end{tabarg}

\medskip
%%%%%SORTIES%%%%
\begin{tabarg}
\multicolumn{3}{l}{\titretab {Sorties}}\\
\hline
$\Theta_0$ & Vecteur, \texttt{float} & Angle solaire z�nithal (degree)\\
$\Phi_{0}$ & Vecteur, \texttt{float} & Angle solaire azimutal (degree)\\
$AMU0$ & Vecteur, \texttt{float} & Cosinus de $\Theta_0$\\
\end{tabarg}


%%%%%FORMULE%%%%
\itemstyle{Formule}

%% Formule de l'operateur : pos_sol_mobile_loa
%% Dans ce fichier, les �ventuels accents sont cod�s
%% en ISO-8859-15 (aka latin9)
%% La ligne suivante donne trois exemples 
%% e accent aigu : �, e accent grave : �, e-dans-l'o : �

Calcul des deux coordonn�es du soleil dans le rep�re �quatorial local :
\[
\begin{cases} \omega &=  15  \frac{\pi}{180} \left(\frac{Time}{3600} + \frac{Lon}{15} - \frac{ET}{60} - 12 \right)\\
\delta &= f\left(2\pi \frac{(JDA-1)}{365}\right)
\end{cases} 
\]
%�quation du temps
$\omega$: angle horaire, l'angle sur le plan �quatorial, 
entre la position du plan m�ridien contenant le soleil et le m�ridien local,
compt�  positivement vers l'Est.
%  	Angle horaire:
%L'une des coordonn�es horaires. Angle di�dre du cercle horaire de la direction et du m�ridien du lieu pris comme
%origine. L'angle horaire est compt� positivement dans le sens r�trograde. 

$\delta$ : d�clinaison, l'angle que forme la direction Terre-Soleil par rapport au 
plan de l'�quateur terrestre, dans le cercle m�ridien contenant le soleil.

$ET$ : l'��quation� du temps, la diff�rence (Temps solaire vrai) - (Temps solaire moyen).

Puis conversion g�om�trique en azimut ($\Phi_0$) et angle solaire z�nithal ($\Theta_0$):
\[
\begin{cases}
Elev &=  \arcsin \left( \sin(Lat) \sin(\delta) + \cos(Lat)\cos(\delta)\cos(\omega) \right)\\
 \Theta_0 &=  90 - Elev\\
\Phi_0 &\text{est tel que } \begin{cases}\sin(\Phi_O)&= \cos(\delta)*\sin(\omega)/\cos(Elev) \\
\cos(\Phi_O)&= \frac{-\cos(Lat)\sin(\delta)+\sin(Lat)\cos(\delta)\cos(\omega)}{\cos(Elev)}
\end{cases} 
\end{cases} 
\]


%%%%%AUTEUR%%%%
\itemstyle{Auteur~:}
Le programme \textsc{fortran} qui fait effectivement les calculs
nous a �t� transmis par G. Brogniez (LOA). L'enrobage
sous forme d'op�rateur a �t� fait par B. Piguet et P. Richard.

%%%%%REFERENCE%%%%
\itemstyle{R\'ef\'erences~:}
G�rard Brogniez, LOA.
